\begin{abstract}
Although large-scale unsupervised language models (LMs) acquire extensive general knowledge and some capacity for reasoning, precisely directing their outputs remains challenging, primarily because their training is entirely unsupervised. Current techniques for enhancing the controllability of LMs often rely on gathering human judgments about the relative merits of model outputs and adapting the base LM to align with these preferences, frequently via reinforcement learning from human feedback (RLHF). However, RLHF procedures are intricate and can be unstable, involving an initial stage where a reward model is trained to mirror human choices, followed by reinforcement learning to adjust the large language model for maximizing the approximated reward, all while attempting to prevent significant deviation from the pre-trained LM. This paper introduces an alternative reward model parameterization for RLHF, enabling the closed-form derivation of the associated optimal policy and reducing the RLHF task to straightforward classification loss optimization. We present \textit{Direct Preference Optimization} (DPO), an algorithm that is robust, efficient, and computationally light, obviating the necessity for sampling from the LM during fine-tuning or for extensive hyperparameter search. Empirical results indicate that DPO tunes LMs to fit human preferences as effectively as, or better than, previous approaches. In particular, DPO fine-tuning outperforms PPO-based RLHF on sentiment control, and attains equivalent or higher quality on summarization and single-turn dialogue, all with much simpler implementation and training procedures.
\end{abstract}

\section{Introduction}
Large-scale unsupervised language models (LMs), when exposed to extensive datasets, demonstrate a range of unexpected abilities~\citep{chowdhery2022palm, brown2020language, touvron2023llama,bubeck2023sparks}. Nevertheless, the human-generated training data encompasses a vast spectrum of intentions, preferences, and expertise levels. Not all of these human attributes are appropriate for AI systems to emulate; for instance, while an AI coding assistant should \textit{recognize} frequent programming errors to effectively correct them, we generally prefer the system to emulate the exceptional, albeit infrequent, examples of expert programming found in the corpus, rather than replicate common mistakes. Likewise, we may desire that a language model be \textit{informed} about a prevalent misconception held by half of the population, but it is undesirable for the model to propagate that falsehood in half of its corresponding answers. Thus, carefully guiding the model's \emph{output and behavioral patterns}—selecting from its broad reservoir of \textit{knowledge and proficiencies}—is essential for producing AI agents that are safe, effective, and amenable to control~\citep{ouyang2022training}. Although current approaches predominantly utilize reinforcement learning (RL) to align LMs with human preferences, this paper demonstrates that the standard RL-based objective adopted in these pipelines is exactly solvable using a straightforward binary cross-entropy objective, which can substantially streamline preference alignment processes.

In general, methods that steer LMs toward desirable behaviors rely on carefully collected datasets of human preference judgments reflecting characteristics that users associate with safety and utility. Such preference-based training supplements the initial unsupervised pre-training conducted over massive corpora of raw text. A direct route for learning human preferences is through supervised fine-tuning on exemplary responses crafted by humans. However, the prevailing and most effective framework is reinforcement learning from human or AI-generated feedback (RLHF/RLAIF; \citep{christiano2017deep,bai2022constitutional}). In RLHF, a reward model is first trained on annotated preference data, after which the LM is fine-tuned via RL to generate outputs that maximize the learned reward function, all while preventing excessive divergence from the pretrained model. Though RLHF yields models with remarkable dialogue and code generation competencies, its methodological complexity is significantly greater than supervised fine-tuning, as it requires maintaining multiple model instances and involves in-the-loop sampling from the policy, ultimately leading to high computational expense.

This work introduces a method to directly tune language models to follow human preferences, foregoing the need for explicit reward modeling or reinforcement learning. We present \textit{{Direct Preference Optimization} (DPO)}, an approach that inherently targets the same objective as RLHF techniques—maximizing reward while regularizing with a KL-divergence penalty—yet is notably easy to implement and train. At a high level, DPO modifies the relative log-probabilities to favor preferred outputs over less favored ones, but further applies a dynamically computed, per-sample weighting to avert the deterioration of model quality that can result from a naïve probability ratio-based objective. DPO, like other algorithms in this space, depends on a theoretical preference framework (e.g., the Bradley-Terry model; \cite{bradley1952rankanalysis}) that quantifies the fit between a reward function and observed human preferences. In contrast to prevailing approaches that first train a reward model with this preference loss and subsequently optimize a policy for the learned reward, DPO leverages a change of variables to reexpress the preference loss directly in terms of the policy. With access to a dataset of human judgments over model responses, DPO can update policies using a straightforward binary cross-entropy loss, thereby attaining the policy that maximizes an implicit reward model tuned to the preference data.

Our principal contribution is Direct Preference Optimization (DPO), a lightweight, reinforcement learning-free algorithm for preference-based training of language models. Empirically, we find that DPO matches or surpasses the effectiveness of leading methods, including those based on PPO for RLHF, across tasks such as sentiment control, text summarization, and dialogue, when applied to models with up to 6B parameters.

\section{Related Work}

As self-supervised language models grow in size, they increasingly demonstrate the capability to solve certain tasks either without any additional training data (zero-shot) \citep{radford2019language} or via a small number of examples provided at inference time (few-shot prompts) \citep{gpt3,megatron,chowdhery2022palm}. Nevertheless, these models’ effectiveness on specific applications and their responsiveness to user instructions are notably enhanced through fine-tuning on corpora comprising task-specific instructions paired with human-crafted responses \citep{mishra-etal-2022-cross,sanh2022multitask,chung2022scaling,thoppilan2022lamda}. Known as `instruction-tuning,' this technique not only enables large language models (LLMs) to adapt to instructions not present in the fine-tuning data, but also substantially boosts their overall practical utility \citep{chung2022scaling}. Although instruction tuning has achieved considerable success, gathering human preference judgments—which assess the quality of model outputs in relative terms—is generally more feasible than obtaining expert demonstrations. As a result, further advancements have been made by refining LLMs using datasets based on such preferences, leading to improvements in translation \citep{kreutzer-etal-2018-reliability}, summarization \citep{stiennon2022learning,ziegler2020finetuning}, narrative generation \citep{ziegler2020finetuning}, and instruction compliance \citep{ouyang2022training,ramamurthy2023is}. These approaches commonly begin with the optimization of a neural reward model that reflects human preferences, often structured using frameworks like the Bradley-Terry model \citep{bradley1952rankanalysis}. Subsequently, the LLM is adapted to maximize the reward predicted by this preference model, employing reinforcement learning strategies such as REINFORCE \citep{williams1992reinforce}, proximal policy optimization (PPO; \cite{schulman2017proximal}), or similar methods \citep{ramamurthy2023is}. A related research direction involves employing instruction-tuned LLMs, which have already been adapted using human feedback, to synthesize additional preference data that target particular criteria such as safety or harmlessness \citep{bai2022constitutional}; in these cases, weak supervision is provided through descriptive rubrics rather than direct human annotations. This trend signifies a merging of research focused on using reinforcement learning to train language models for diverse goals~\citep{Ranzato2015SequenceLT,paulus2018a,wu2018learning} and efforts aimed at general frameworks for extracting knowledge from human preferences \citep{christiano2017deep,kupcsik2018learning}. In spite of the advantages of employing relative human judgments, scaling reinforcement learning-based fine-tuning to large language models presents significant implementation obstacles; the current work introduces a theoretically grounded methodology for optimizing with relative preference signals without the reliance on RL.

Research on learning policies from preferences—outside the domain of language—has been explored in both bandit and reinforcement learning frameworks, resulting in multiple methodologies. In contextual bandit scenarios, preferences or action rankings can substitute for rewards; this is encapsulated by the contextual dueling bandit (CDB) paradigm \citep{yue2012karmed,dudik2015contextual}. When explicit reward signals are missing, the theoretical treatment of CDBs replaces the concept of an optimal policy with that of a \textit{von Neumann winner}: a policy whose expected win probability is at least 50\% against \textit{every} alternative policy \citep{dudik2015contextual}. A crucial difference between CDBs and human preference learning is that in CDB, preference feedback is provided interactively (online), whereas learning from human judgments typically involves a static dataset of pre-collected, preference-labeled action pairs \citep{yan2022human}. Likewise, in \textit{preference-based RL} (PbRL), learning is guided by binary comparisons generated via an unknown `scoring' function, not from reward signals \citep{BusaFekete2014,ruiz2023dueling}. PbRL algorithms often rely on reconstructing this underlying scoring function—essentially inferring a reward model—prior to or as part of policy optimization, and may accommodate off-policy data \citep{jain2013learning,BusaFekete2014,christiano2017deep,sadigh2017active,kupcsik2018learning}. Contrary to these two-step procedures, our approach unifies policy learning into a single stage, directly aligning policy outputs to preference information.

\section{Preliminaries}\label{section:prelims}

We provide an overview of the RLHF workflow detailed by \citeauthor{ziegler2020finetuning} (and subsequent works: \citep{stiennon2022learning, bai2022training, ouyang2022training}), which is generally divided into three main stages: (1) supervised fine-tuning (SFT); (2) preference elicitation and reward modeling; and (3) policy optimization via RL.

\textbf{SFT}: The RLHF procedure typically initiates with supervised fine-tuning of a pre-trained language model (LM) using high-quality datasets tailored to the downstream objective(s) (such as dialogue or summarization). This yields the model $\pi^\text{SFT}$.

\textbf{Reward Modeling Stage}: In this second step, prompts $x$ are fed to the SFT-trained model to sample answer pairs $(y_1, y_2)\sim \pi^\text{SFT}(y \mid x)$. Human evaluators are then asked to choose their preferred completion, expressed as $y_w\succ y_l \mid x$, where $y_w$ represents the favored answer, and $y_l$ the less favored one from the pair $(y_1, y_2)$. The presumption is that such preferences are governed by a latent reward function $r^*(y, x)$, inaccessible to the learner. Preferences may be modeled through a variety of frameworks, with the Bradley-Terry (BT) model \cite{bradley1952rankanalysis} being frequently used (though the framework can accommodate more comprehensive ranking models such as Plackett-Luce \citep{plackett1975analysis, luce2012individual}, especially if multiple answers are ranked). The BT approach posits the human preference probabilities $p^*$ can be formally captured by:

Given a static comparison dataset $\mathcal{D} = \bigl\{x^{(i)}, y_w^{(i)}, y_l^{(i)}\bigr\}_{i=1}^N$ drawn from $p^*$, a reward model parameterized by $r_\phi(x, y)$ can be fit via maximum likelihood estimation. Casting this as a binary classification scenario, the negative log-likelihood loss is formulated as follows:

\begin{equation}\label{eq:reward_model}
    \mathcal{L}_R(r_{\phi}, \mathcal{D}) = -\mathbb{E}_{(x, y_w, y_l)\sim \mathcal{D}}\bigl[\log \sigma(r_{\phi}(x, y_w)- r_{\phi}(x, y_l))\bigr]
\end{equation}

with $\sigma$ representing the logistic sigmoid function. Within the context of language models, $r_\phi(x, y)$ is commonly initialized using the SFT model $\pi^\text{SFT}(y \mid x)$, with an additional linear transformation on top of the final transformer layer to yield a scalar reward prediction \cite{ziegler2020finetuning}. To minimize reward estimate variance, previous studies standardize reward values so that $\mathbb{E}_{x,y\sim \mathcal{D}}\left[r_\phi(x, y)\right] = 0$ for every $x$.

\textbf{RL Fine-Tuning Phase}: When proceeding to reinforcement learning, the learned reward model provides feedback for further training the language model. Following earlier methodologies~\citep{jaques2017sequence, jaques2020human}, the policy optimization objective is given by

\begin{equation}\label{eq:RL}
\max_{\pi_{\theta}}  \mathbb{E}_{x\sim \mathcal{D}, y\sim \pi_{\theta}(y \mid x)}\bigl[r_{\phi}(x, y)\bigr] - \beta\mathbb{D}_{\textrm{KL}}\bigl[\pi_{\theta}(y\mid x)\mid \mid \pi_\text{ref}(y\mid x)\bigr],
\end{equation}

where $\beta$ governs the allowable divergence from a reference policy $\pi_\text{ref}$—typically, this is taken to be the initial SFT model $\pi^\text{SFT}$. The policy $\pi_\theta$ is also usually initialized from $\pi^\text{SFT}$. This regularization term is essential; it discourages excessive departures from the region of distributional support where the reward model is reliable, helps preserve output diversity, and mitigates mode collapse onto a small set of high-reward responses. Because language generation is inherently discrete and non-differentiable, this policy optimization problem is commonly approached via reinforcement learning. The usual procedure \citep{ziegler2020finetuning, stiennon2022learning, bai2022training, ouyang2022training} entails defining a reward function as ${r(x, y) = r_{\phi}(x, y) -\beta (\log \pi_{\theta}(y\mid x) - \log \pi_\text{ref}(y\mid x))}$ and applying PPO \cite{schulman2017proximal} to maximize the resulting objective.

\section{Direct Preference Optimization}\label{sec:DPO}

Faced with the practical challenges associated with applying standard reinforcement learning approaches to large models, particularly during language model fine-tuning, our objective is to propose a more streamlined method for policy optimization grounded in preference data. Departing from the typical RLHF paradigm—where a reward is learned and then optimized with RL—we instead employ a reward model parametrization that admits a closed-form solution for the associated optimal policy, thereby eliminating the need for an explicit RL training loop. Central to our proposal is the observation that a direct analytic mapping can be established between reward models and their optimal policies. This mapping enables us to recast loss functions defined over reward models in terms of policy loss functions. In this framework, an explicit standalone reward network is not necessary, yet we remain consistent with established models of human preference, such as the Bradley-Terry formulation. Under this approach, the policy network serves jointly as both the language generator and the (implicitly defined) reward estimator.

\textbf{Derivation of the DPO Objective.} We begin from the general RL framework as established in earlier works, i.e., Eq.~\ref{eq:RL}, considering an arbitrary reward function $r$. Building upon prior literature~\citep{peters2007reinforcement, peng2019advantage, korbak2022reinforcement, go2023aligning}, it follows directly that the policy optimizing the KL-regularized reward maximization objective in Eq.~\ref{eq:RL} admits the solution:

\begin{equation}\label{eq:op_policy}
    \pi_r(y\mid x) = \frac{1}{Z(x)}\pi_\text{ref}(y\mid x)\exp\left(\frac{1}{\beta}r(x, y)\right),
\end{equation}

where the normalization constant $Z(x) = \sum_{y}\pi_\text{ref}(y\mid x)\exp\left(\frac{1}{\beta}r(x, y)\right)$ is known as the partition function (refer to Appendix~\ref{app:derivation1} for details). Even if one adopts the maximum likelihood estimate $r_{\phi}$ of the actual reward $r^*$, direct computation of $Z(x)$ remains computationally demanding~\citep{korbak2022reinforcement, go2023aligning}, impeding practical use of this formulation. Nonetheless, it is possible to reformulate Eq.~\ref{eq:op_policy} so that the reward is written in terms of the optimal policy $\pi_r$, the baseline $\pi_\text{ref}$, and the partition function $Z(\cdot)$. Specifically, taking the logarithm of both sides of Eq.~\ref{eq:op_policy} and rearranging yields:

\begin{equation}\label{eq:main_eq}
    r(x,y) = \beta \log \frac{\pi_r(y\mid x)}{\pi_\text{ref}(y\mid x)} + \beta \log Z(x).
\end{equation}

Applying this transformation to the true reward $r^*$ and its associated optimal policy $\pi^*$, we leverage the fact that the Bradley-Terry model is influenced only by reward differences between two outputs, i.e., ${p^*(y_1 \succ y_2 \mid x) = \sigma(r^*(x, y_1) - r^*(x, y_2))}$. Substituting Eq.~\ref{eq:main_eq} (with $r^*$ and $\pi^*$) into the preference expression Eq.~\ref{eq:bradley-terry}, the partition function $Z(x)$ cancels, leaving a dependence solely on the optimal and reference policies. Hence, under the Bradley-Terry model, the optimal RLHF policy $\pi^*$ satisfies:

\begin{equation}\label{eq:objective}
    p^*(y_1\succ y_2 \mid x) = \frac{1}{1 + \exp\left(\beta \log \frac{\pi^*(y_2\mid x)}{\pi_\text{ref}(y_2\mid x)} - \beta \log \frac{\pi^*(y_1\mid x)}{\pi_\text{ref}(y_1\mid x)}\right)}
\end{equation}

A step-by-step derivation is provided in Appendix~\ref{app:derivation2}. Although Eq.~\ref{eq:objective} is based on the Bradley-Terry formulation, analogous developments can be made for the more general Plackett-Luce family of models~\citep{plackett1975analysis, luce2012individual}, as discussed in Appendix~\ref{app:plackett_luce_models}.

Given this formulation, where the human preference likelihood is written directly in terms of the policy rather than an explicit reward model, it becomes natural to define a maximum likelihood estimation procedure for a parameterized policy $\pi_\theta$. This is structurally similar to standard reward model objectives (see Eq.~\ref{eq:reward_model}), but is now re-cast as a policy learning problem:

\begin{equation}\label{eq:optimum_model}
    \mathcal{L}_\text{DPO}(\pi_{\theta}; \pi_\text{ref}) = -\mathbb{E}_{(x, y_w, y_l)\sim \mathcal{D}}\left[\log \sigma \left(\beta \log \frac{\pi_{\theta}(y_w\mid x)}{\pi_\text{ref}(y_w\mid x)} - \beta \log \frac{\pi_{\

In this approach, we learn an implicit reward function via an alternative form of parameterization, for which the optimal policy coincides with $\pi_\theta$. Additionally, our method aligns with fitting a reparametrized Bradley-Terry model, and thus inherits desirable theoretical characteristics, such as consistency provided the preference data distribution meets certain assumptions \cite{bong2022generalized}. A more detailed exploration of the theoretical aspects of DPO and its connection to related methodologies appears in Section~\ref{sec:theory}.

\textbf{How does the DPO update function?} To gain mechanistic insight into DPO, it is instructive to consider the gradient of its objective, $\mathcal{L}_\text{DPO}$. The gradient with respect to the parameters $\theta$ is expressed as:

\begin{multline*}\label{eq:gradient}
    \nabla_\theta \mathcal{L}_\text{DPO}(\pi_\theta;\pi_\text{ref}) = \\ -\beta\mathbb{E}_{(x, y_w, y_l) \sim \mathcal{D}} \bigg[\underbrace{\sigma(\hat{r}_\theta(x, y_l) - \hat{r}_\theta (x, y_w))}_\text{greater emphasis when reward assessment errs}\bigg[\underbrace{\nabla_\theta\log \pi(y_w \mid x)}_\text{boost probability of $y_w$} - \underbrace{\nabla_\theta\log\pi(y_l \mid x)}_\text{reduce probability of $y_l$}\bigg]\bigg],
\end{multline*}

where $\hat{r}_\theta(x, y) = \beta \log \frac{\pi_\theta(y \mid x)}{\pi_\text{ref}(y \mid x)}$ designates the reward implicitly specified by both the target language model $\pi_\theta$ and the reference model $\pi_\text{ref}$ (see Section~\ref{sec:theory} for elaboration). From an intuitive perspective, the gradient $\mathcal{L}_\text{DPO}$ encourages higher likelihood for the preferred completions $y_w$ while discouraging the less-preferred $y_l$. Crucially, the impact of each example is scaled in accordance with how much the implicit reward model $\hat{r}_\theta$ overestimates the disfavored outputs, adjusted by the factor $\beta$; that is, the update intensity reflects the magnitude of error in ranking the completions, and incorporates the influence of the KL regularization. Empirical results underscore the necessity of this weighting: omitting the weighting factor can lead to severe model collapse, as shown in Appendix Table~\ref{tab:unlikelihood_generations}.

\textbf{DPO overview.} The standard DPO workflow consists of the following steps: First, for each prompt $x$, generate candidate completions $y_1, y_2 \sim \pi_\text{ref}(\cdot \mid x)$ and assign preference labels based on human judgments to assemble an offline preference dataset $\mathcal{D} = \{x^{(i)}, y_w^{(i)}, y_l^{(i)}\}_{i=1}^N$. Second, train the language model $\pi_\theta$ by minimizing $\mathcal{L}_\text{DPO}$ given the reference model $\pi_\text{ref}$, the dataset $\mathcal{D}$, and the chosen parameter $\beta$. In practical settings, instead of creating new samples and collecting fresh human preference data, practitioners typically leverage pre-existing publicly available preference datasets. Because these datasets are often produced via $\pi^\text{SFT}$, the reference policy is typically initialized as $\pi_\text{ref} = \pi^\text{SFT}$ when possible. If $\pi^\text{SFT}$ is not accessible, $\pi_\text{ref}$ is set by maximizing the likelihood over preferred completions ${(x, y_w)}$, specifically, ${\pi_\text{ref} = \argmax_{\pi}\mathbb{E}_{x, y_w \sim \mathcal{D}}\left[\log \pi(y_w \mid x)\right]}$. This strategy addresses the distribution mismatch between the ideal but unavailable reference distribution and the operational $\pi_\text{ref}$ in DPO. Comprehensive information regarding implementation specifics and the choice of hyperparameters is provided in Appendix~\ref{app:implementation}.

\section{Theoretical Analysis of DPO}
This section develops a theoretical understanding of the DPO approach, interpreting its mechanisms, supplying formal justifications, and contrasting its benefits with the challenges encountered by actor-critic-based RLHF algorithms (e.g., PPO~\cite{schulman2017proximal}).

\label{sec:theory}

\subsection{Your Language Model Is Secretly a Reward Model}
DPO eliminates the need for constructing an explicit reward model and bypasses reinforcement learning, instead utilizing a single maximum likelihood objective. The loss function in Eq. \ref{eq:main_eq} matches the Bradley-Terry model structure when the rewards are parameterized as $r^*(x, y) = \beta \log\frac{\pi^*_\theta(y \mid x)}{\pi_\text{ref}(y \mid x)}$; consequently, optimizing the parametric family $\pi_\theta$ corresponds to optimizing the reward model as in Eq. \ref{eq:reward_model}, but under a reparameterized framework. This section lays the foundation for this reparameterization, establishes that it does not restrict the space of attainable reward models, and shows it can exactly recover the optimal policy. We begin by introducing an equivalence relation on reward functions.

\begin{definition}
We define two reward functions $r(x, y)$ and $r'(x, y)$ as equivalent if and only if ${r(x, y)-r'(x, y) = f(x)}$ holds for some function $f$.
\end{definition}

This relation evidently satisfies the properties of an equivalence relation, thereby dividing the collection of reward functions into equivalence classes. This allows us to state the following lemmas:

\begin{lemma}\label{lemma:same_prefrence}
Within the framework of Plackett-Luce models, and specifically the Bradley-Terry model, any pair of reward functions belonging to the same equivalence class will yield identical preference distributions.
\end{lemma}

\begin{lemma}\label{lemma:same_policy}
    For the constrained RL problem, all reward functions that are members of the same equivalence class result in the same optimal policy.
\end{lemma}

The arguments supporting these statements are straightforward, and we postpone them to Appendix \ref{app:lemma1}. The first lemma highlights the well-established issue of model under-specification in the Plackett-Luce family \cite{plackett1975analysis}. Owing to this lack of identifiability, it is typically necessary to adopt further constraints in order to obtain meaningful maximum likelihood estimates as in Eq. \ref{eq:reward_model} \cite{bong2022generalized}. The second lemma establishes that any reward function within a given class will produce the same optimal policy; consequently, for the purpose of optimization, recovering any member of the optimal equivalence class suffices. The theorem below, whose proof appears in Appendix~\ref{app:thm1}, formalizes this observation:

\begin{theorem}\label{thm:main}
    Provided mild regularity conditions, every equivalence class of reward functions compatible with Plackett-Luce (and, in particular, Bradley-Terry) models admits a representation of the form ${r(x, y) = \beta \log \frac{\pi(y\mid x)}{\pi_\text{ref}(y\mid x)}}$ for some model $\pi(y\mid x)$ and any fixed reference model $\pi_\text{ref}(y \mid x)$.
\end{theorem}

\begin{sproof}
    Take an arbitrary reward function $r(x, y)$, which determines a corresponding optimal model $\pi_r(y \mid x)$ according to Eq. \ref{eq:op_policy}. We aim to show that a reward function in the same equivalence class as $r$ may be expressed using the above parameterization. Define the projection operator $f$ by
\begin{equation}
    f(r; \pi_\text{ref}, \beta)(x, y) = r(x, y) - \beta\log\sum_{y}\pi_\text{ref}(y\mid x)\exp\left(\frac{1}{\beta}r(x, y)\right)
\end{equation}
This operator normalizes the reward by subtracting the logarithm of the partition function (associated with $\pi_r$), and since this term is only dependent on $x$, $f(r; \pi_\text{ref}, \beta)(x, y)$ remains within the equivalence class of $r(x, y)$. Now, substituting $r$ with the right-hand side of Eq.~\ref{eq:main_eq}, which is valid for any $r$, yields $f(r; \pi_\text{ref}, \beta)(x, y) = \beta \log \frac{\pi_r(y\mid x)}{\pi_\text{ref}(y\mid x)}$. Therefore, $f$ produces a representative of the equivalence class of $r$ in the desired parameterized form, ensuring no loss of generality in the reward model by employing this reparameterization.
\end{sproof}

An alternative interpretation of Theorem~\ref{thm:main} is that it precisely identifies, within each equivalence class, the specific reward function selected by the DPO reparameterization: namely, the reward function satisfying

\begin{equation}\label{eq:lag_p}
     \sum_{y}\underbrace{\pi_\text{ref}(y\mid x)\exp\left(\frac{1}{\beta}r(x, y)\right)}_{=\pi(y\mid x)\text{, as per Thm.~\ref{thm:main} reparam.}} = 1,
\end{equation}

which ensures that $\pi(y\mid x)$ forms a legitimate probability distribution—that is, all probabilities are non-negative and collectively sum to unity. Moreover, in light of Eq.~\ref{eq:op_policy}, Eq.~\ref{eq:lag_p} can be recognized as corresponding to the partition function associated with the optimal policy defined by the reward function $r(x, y)$. A crucial conceptual advancement in the DPO algorithm is its introduction of specific constraints on the typically under-determined Plackett-Luce (and, in particular, Bradley-Terry) family of preference models. These constraints retain the full expressive power for representing reward models while, at the same time, rendering the optimal policy described by Eq.~\ref{eq:op_policy} explicitly computable for any prompt $x$.

\subsection{Instability of Actor-Critic Algorithms} Our framework further facilitates the identification of instability issues in standard actor-critic algorithms employed in RLHF pipelines, such as PPO. Specifically, we consider the RL fine-tuning phase as detailed in Section \ref{section:prelims}. By relating this setting to the control as inference paradigm \cite{levine2018reinforcement} for the constrained reinforcement learning problem described by \ref{eq:RL}, we assume a parameterized model $\pi_{\theta}(y\mid x)$ and seek to minimize $\mathbb{D}_{\text{KL}}[\pi_{\theta}(y|x) \mid \mid \pi^*(y\mid x)]$, where $\pi^*$, defined via Eq.~\ref{eq:optimum_model}, represents the optimal policy dictated by the reward $r_{\phi}(y, x)$. Through some algebraic manipulations, we arrive at the following objective:

\begin{equation}\label{eq:AC}
    \max_{\pi_{\theta}}\mathbb{E}_{\pi_{\theta}(y\mid x)}\bigg[\underbrace{r_{\phi}(x, y) -\beta\log\sum_{y}\pi_\text{ref}(y\mid x)\exp\left(\frac{1}{\beta}r_{\phi}(x, y)\right)}_{f(r_{\phi}, \pi_\text{ref}, \beta)} - \underbrace{\beta\log\frac{\pi_{\theta}(y\mid x)}{\pi_\text{ref}(y\mid x)}}_{\text{KL}}\bigg]
\end{equation}

The objective considered here aligns with those optimized in previous studies 
\citep{ziegler2020finetuning, stiennon2022learning, bai2022training, ouyang2022training}, which employ the DPO-equivalent reward structure for the reward function $r_{\phi}$. Under this formulation, the normalization component in $f(r_{\phi}, \pi_\text{ref}, \beta)$ may be viewed as the soft value function corresponding to the reference policy $\pi_\text{ref}$. Although this term does not change the optimal policy, omitting it can significantly increase the variance of policy gradients, thereby destabilizing training. To address this, a value function can be learned to model the normalization term, but optimizing such a function presents its own challenges. Alternatively, previous methods have utilized a baseline derived from human completions to normalize rewards, effectively providing a single-sample Monte Carlo estimate of the normalization constant. In contrast, the DPO reparameterization framework results in a reward function that inherently eliminates the need for such baselines.

\section{Experiments}
We now turn to an empirical assessment of DPO’s capacity to train policies directly from preference data. To begin, within a controlled text generation scenario, we examine how effectively DPO balances reward maximization and KL-divergence minimization with respect to the reference policy, benchmarking it against standard preference learning algorithms like PPO. We then extend our analysis to encompass larger-scale models and more challenging RLHF applications, such as summarization and dialog tasks. Our experiments demonstrate that DPO, with minimal hyperparameter adjustment, typically matches or outperforms robust baselines—including RLHF with PPO and the selection of the top trajectory from $N$ samples using a learned reward function. Prior to discussing empirical findings, we detail the experimental methodology; further information can be found in Appendix~\ref{app:exp_details}.

\textbf{Tasks.} The scope of our experimental study spans three distinct open-ended text generation problems. In all settings, models are tasked with learning a policy based on a collection of preference data $\mathcal{D}=\bigl\{x^{(i)}, y_w^{(i)}, y_l^{(i)}\bigr\}_{i=1}^N$. For the \textbf{controlled sentiment generation} task, $x$ represents a snippet from a movie review in the IMDb dataset \cite{maas-EtAl:2011:ACL-HLT2011}, and the objective is to generate a continuation $y$ with positive sentiment. To facilitate objective assessment, preference pairs for this task are generated by applying a pre-trained sentiment classifier, ensuring $p(\text{positive}\mid x,y_w)>p(\text{positive}\mid x,y_l)$. In this context, GPT-2-large is fine-tuned using the training subset of the IMDB reviews to obtain an SFT model (implementation details in App~\ref{app:sentiment_details}). For \textbf{summarization}, the input $x$ consists of a Reddit forum post, and the policy must compose a summary $y$ that distills the main ideas from the post. This setup employs the Reddit TL;DR summarization corpus \citep{volske-etal-2017-tl}, supplemented by human preference annotations from \citeauthor{stiennon2022learning}. The SFT model is obtained via fine-tuning on summaries written by humans\footnote{\url{https://huggingface.co/CarperAI/openai_summarize_tldr_sft}}, with RLHF training facilitated through the TRLX library \citep{leandro_von_werra_2023_7790115}. The preferences are sourced from the human preference dataset curated by \citeauthor{stiennon2022learning}, based on completions from a similarly-trained, but different, SFT model. The third scenario, \textbf{single-turn dialogue}, involves $x$ as a single user prompt that may range from scientific inquiries to personal advice requests. Here, the aim is for the policy to craft a response $y$ that is both useful and engaging; the Anthropic Helpful and Harmless dialogue dataset \citep{bai2022training} serves as the benchmark, containing 170k exchanges between human users and an automated assistant. Each sample ends with two model-generated responses and a human annotation of the preferred choice. Since no ready-made SFT model exists for this data, we train an SFT model from scratch by fine-tuning a pre-existing language model solely on the human-preferred responses.

\textbf{Evaluation.} We employ two distinct strategies to evaluate our methods. In the context of controlled sentiment generation, where the true reward function (provided by a sentiment classifier) is accessible, we assess each algorithm by tracing the Pareto frontier of obtained reward versus KL-divergence from the original policy. This direct comparison is feasible thanks to the availability of ground-truth rewards. In contrast, in practical applications where the reward function is unobservable, we rely on the \textit{win rate} relative to a baseline policy for evaluation. For this, GPT-4 serves as an automated surrogate for human judgment, assessing summary quality in the summarization setting and response helpfulness in single-turn dialogue. Baselines are defined by the reference summaries within the test split for summarization, and by the human-preferred responses in the test split for dialogue. While some literature supports the reliability of language models as evaluators over standard metrics \citep{Chen2023ExploringTU}, we validate our use of GPT-4 by conducting an accompanying human assessment study, discussed in Sec.~\ref{sec:human-judgments}. The results demonstrate a strong alignment between GPT-4’s evaluations and those of human raters; agreement between humans and GPT-4 matches or exceeds agreement levels among different human annotators.

\textbf{Methods.} Beyond DPO, we benchmark several established methods for preference-aligned language model training. For the summarization task, we utilize zero-shot prompting with \textbf{GPT-J} \citep{gpt-j}; for the dialogue task, we employ two-shot prompting using \textbf{Pythia-2.8B} \citep{biderman2023pythia}. Further, we assess a \textbf{SFT} model as well as the \textbf{Preferred-FT} variant, which is trained via supervised fine-tuning on the selected output $y_w$—this selection comes from SFT in sentiment/summarization and from a general LM in dialogue. Another approach under consideration is \textbf{Unlikelihood}~\citep{welleck2019neural}, which updates the policy to increase the likelihood of $y_w$ while simultaneously decreasing the probability of $y_l$; an adjustable parameter $\alpha\in[0,1]$ scales the penalty term. We further include \textbf{PPO} \citep{schulman2017proximal}, where the reward model is learned from preference feedback, and an oracle variant, \textbf{PPO-GT}, which has access to the exact reward function in the sentiment task. Our sentiment experiments leverage two versions of PPO-GT: a standard implementation \cite{leandro_von_werra_2023_7790115}, and a revised form featuring reward normalization and refined hyperparameter tuning for enhanced results. These enhancements are also applied when running PPO with learned rewards. As a strong baseline, we include the \textbf{Best of $N$} strategy: here, $N$ responses are generated from the SFT model (or Preferred-FT for dialogue) and the candidate with the highest score—according to a preference-trained reward model—is selected. While this technique often performs very well, it requires sampling $N$ outputs per test prompt, rendering it computationally expensive and thus impractical for larger values of $N$.

\subsection{How well can DPO optimize the RLHF objective?}

Typical RLHF algorithms are designed to maximize rewards while constraining the KL divergence from a reference policy, thereby promoting reward-seeking behavior without allowing the policy to diverge excessively. Consequently, in evaluating such algorithms, both the attained reward and the corresponding KL divergence must be considered; a policy that secures marginally greater reward at the expense of significantly increased KL divergence may be suboptimal. In Figure~\ref{fig:frontier-tldr-main}, we present the reward versus KL tradeoff for different algorithms under the sentiment task. For each method, we perform several training sessions, adjusting the policy conservativeness parameter for each session (with target KL $\in\{3,6,9,12\}$ for PPO, $\beta \in \{0.05,0.1,1,5\}$, $\alpha\in\{0.05,0.1,0.5,1\}$ for unlikelihood training, and varying random seeds for preferred-FT). This procedure yields a total of 22 distinct runs. After every 100 training iterations, up to convergence, we evaluate all policies using a set of held-out test prompts, measuring the average reward with respect to the true reward function, as well as the mean sequence-level KL\footnote{That is, the aggregate sum of the KL-divergence computed at each timestep.} relative to the reference policy $\text{KL}\left(\pi\mid \mid \pi_\text{ref}\right)$. The results reveal that DPO constructs the most favorable tradeoff curve, simultaneously attaining the highest rewards and keeping KL divergence low. This outcome stands out for several reasons. Notably, DPO and PPO target the same optimization objective, yet DPO consistently outperforms PPO in terms of the reward-KL tradeoff, fully dominating it. Moreover, DPO secures a superior tradeoff curve even in comparison to PPO equipped with access to the actual reward values (PPO-GT).

\subsection{Can DPO scale to real preference datasets?}
\label{sec:dpo-real-datasets}
We proceed to investigate how well DPO performs during fine-tuning on real-world preference tasks in summarization and single-turn dialogue settings. In the case of summarization, established automated metrics like ROUGE often exhibit weak alignment with actual human judgments~\citep{stiennon2022learning}, motivating the use of methods such as PPO fine-tuned on human preferences to enhance summary quality. Our evaluation involves generating completions from the test partition of the TL;DR summarization dataset and calculating the mean win rate of each method relative to the reference summaries in this test set. For all compared techniques, completions are drawn at various sampling temperatures from 0.0 to 1.0, with corresponding win rates depicted in Figure~\ref{fig:frontier-tldr-main} (right). DPO, PPO, and Preferred-FT are all initialized from the identical GPT-J SFT checkpoint\footnote{\url{https://huggingface.co/CarperAI/openai_summarize_tldr_sft}}. At temperature 0.0, DPO achieves a win rate close to 61\%, outperforming PPO, which attains approximately 57\% at its own optimal temperature (also 0.0). DPO additionally surpasses the strongest performance of the best of $N$ baseline. It is important to point out that the DPO $\beta$ hyperparameter was not systematically optimized for these experiments, so the observed results may fall short of DPO's true capabilities. Furthermore, DPO demonstrates considerably greater resilience to changes in sampling temperature than PPO, whose win rate declines toward the baseline GPT-J model as temperature increases. The Preferred-FT approach, by contrast, shows little to no gain over standard SFT. Section~\ref{sec:human-judgments} further details direct human comparisons between DPO and PPO; here, DPO samples at a temperature of 0.25 were selected as preferable in 58\% of cases when compared to PPO samples generated at a temperature of 0.

For single-turn dialogue tasks, we assess the various approaches using the portion of the Anthropic HH test split \citep{bai2022training} that contains exactly one exchange between a human and the assistant. The GPT-4-based evaluations treat the human-preferred responses in the test set as ground truth to calculate the win rate for each technique. Since no standard SFT model exists for this scenario, our procedure begins with the pre-trained Pythia-2.8B. We first apply Preferred-FT to finetune a reference model on the selected completions so that the output remains consistent with the reference model’s data distribution, after which we conduct further training with DPO. We also benchmark against the best outcome among 128 Preferred-FT completions (having determined that the Best of $N$ metric saturates at $N=128$ for this problem; refer to Appendix Figure~\ref{fig:best-of-n}), as well as against a 2-shot prompt variant of the base Pythia-2.8B model. Across optimal temperature settings for each strategy, DPO matches or exceeds the leading alternatives. Additionally, we test an RLHF system, trained via PPO on the Anthropic HH dataset \footnote{\url{https://huggingface.co/reciprocate/ppo_hh_pythia-6B}} following a widely recognized implementation \footnote{\url{https://github.com/CarperAI/trlx/tree/main/examples/hh}}, but we are unable to find prompting or sampling temperature configurations that outperform the original Pythia-2.8B baseline. Given these outcomes, as well as previous results from the TL;DR task and the alignment in reward function objectives, we regard Best of 128 as an approximate indicator of PPO-level effectiveness. In summary, DPO stands out as the only method in this study that both efficiently leverages computation and surpasses the human-preferred completions of the Anthropic HH set, yielding results that match or exceed the resource-intensive Best of 128 benchmark. Notably, Figure~\ref{fig:dialogue-main} demonstrates that DPO attains peak performance within a small number of training iterations.

\subsection{Generalization to a new input distribution}

To more thoroughly assess how PPO and DPO handle distributional shifts, we analyze the policies learned by both methods in our Reddit TL;DR summarization task by applying them to an out-of-domain dataset: the test split of the CNN/DailyMail news articles \citep{nallapati-etal-2016-abstractive}. We employ the optimal sampling temperatures established from TL;DR (0 and 0.25) for this evaluation. Table~\ref{tab:ood} summarizes these findings. To gauge performance, we compute the win rate using GPT-4 against the gold-standard summaries for each dataset, leveraging the same GPT-4 (C) prompt from the TL;DR evaluation, with the term ``forum post'' swapped out for ``news article''. Remarkably, DPO maintains a notable lead over PPO in this new context. These observations offer preliminary support that DPO-trained policies achieve generalization abilities on par with those of PPO, despite DPO not leveraging the additional unlabeled Reddit TL;DR prompts utilized by PPO.

\subsection{Validating GPT-4 judgments with human judgments}
\label{sec:human-judgments}

To assess the trustworthiness of GPT-4's scoring, we run a human evaluation based on the TL;DR summarization outcomes and two distinct GPT-4 prompts. The \textbf{GPT-4 (S)} (simple) prompt asks only which summary better encapsulates the key content of a post, whereas the \textbf{GPT-4 (C)} (concise) prompt incorporates an explicit preference for conciseness; this addition is motivated by observations that GPT-4, when given the \textbf{GPT-4 (S)} prompt, favors summaries that are longer and more redundant than those preferred by humans. Full prompt wording can be found in Appendix~\ref{app:prompts}. We conduct three sets of pairwise model comparisons that span a range of summary qualities: the top-performing DPO at temperature 0.25, the bottom-performing PPO at temperature 1.0, and an intermediate SFT at temperature 0.25, each pitted against PPO with its best-performing temperature (greedy decoding). Results indicate that, for both prompts, agreement between GPT-4 and human raters is on par with inter-human agreement rates, which suggests that GPT-4 may serve as a useful stand-in for human judgment (noting that we gathered multiple human annotations only for the DPO and PPO-1 comparisons, due to a limited pool of raters). In summary, the \textbf{GPT-4 (C)} prompt typically yields win rates more aligned with human preference, and therefore we base the principal results in Section~\ref{sec:dpo-real-datasets} on this prompt. For more information about the human rating process, including the web interface and rater list, consult Appendix~\ref{app:human-study}.

\section{Discussion}
Preference-based learning constitutes an effective and adaptable approach for developing robust and value-aligned language models. In this work, we presented DPO, a straightforward methodology that facilitates preference-based language model training without relying on reinforcement learning. Unlike prior techniques that reformulate the preference learning challenge into an RL paradigm to leverage established RL algorithms, DPO establishes an explicit relationship between language model policies and reward structures. This mapping allows for language model optimization in line with human preferences using a direct cross-entropy loss function, thereby dispensing with reinforcement learning and avoiding a loss of generality. With almost no hyperparameter adjustment required, DPO matches or surpasses the performance of prevailing RLHF techniques, such as those built on PPO, thus significantly lowering the barriers for leveraging human preferences in language model training.

\textbf{Limitations \& Future Work.} Several avenues remain for deeper investigation. To what extent can DPO-trained policies transfer to out-of-distribution data compared to those trained using explicit reward functions? Preliminary evidence indicates that DPO policies generalize comparably to those optimized with PPO, though a more exhaustive evaluation is warranted. An open question is whether integrating self-labeling strategies within DPO could similarly capitalize on unlabeled prompt data. Furthermore, understanding the manifestation of reward over-optimization within the direct preference optimization framework deserves careful analysis—such as whether the modest decline observed in Figure~\ref{fig:dialogue-main}-right exemplifies this phenomenon. While our experiments focus on models of up to 6B parameters, scaling DPO to the largest contemporary language models remains an important and promising research direction. On the subject of evaluation, our findings suggest that GPT-4-determined win rates can vary depending on prompt phrasing; further research is needed to determine optimal strategies for eliciting reliable judgments from automated evaluators. Beyond language modeling from human preferences, DPO also offers considerable potential for application in training generative models across diverse modalities.

\end{document}